\section*{Chapter \ref{ch:b1:electronic-effects} --- Electronic Effects and Reactive Intermediates}

\begin{solution}{exo:b1:electronic-effects:1}
2-Methylbutane, \ce{CH3-CH(CH3)-CH2-CH3}: the three \ce{CH3} are primary,
the \ce{CH2} secondary, the CH tertiary. 2-Bromo-2-methylpropane: the three
\ce{CH3} are primary, the central carbon tertiary. Of the four bromides
\ce{C4H9Br} (1-bromobutane, 2-bromobutane, 1-bromo-2-methylpropane,
2-bromo-2-methylpropane), the last gives the tertiary \omterm{def:b1:electronic-effects:intermediates}{carbocation}, the
most stable.
\end{solution}

\begin{solution}{exo:b1:electronic-effects:2}
$10^{4.76 - 2.87} = 10^{1.89} = 78$; $10^{4.76 - 0.52} = 10^{4.24} =
\num{1.7e4}$.
\end{solution}

\begin{solution}{exo:b1:electronic-effects:3}
Ethanoate: the C=O and \ce{C-O^-} exchanged between the two oxygens.
4-Nitrophenoxide: the oxygen \omterm{def:b1:lewis-resonance-vsepr:lewis}{lone pair} forms a C=O, the ring double bonds
shift, the carbon bearing \ce{NO2} forms a C=N bond, and an N=O pair goes
to oxygen: a structure with C=O at one end of the ring, C=N at the other
and both nitro oxygens negative (a quinoid structure). Allyl cation:
\ce{CH2=CH-CH2+} $\leftrightarrow$ \ce{+CH2-CH=CH2}.
\end{solution}

\begin{solution}{exo:b1:electronic-effects:4}
Second arrow: from the O--H bond of water to its oxygen; products \ce{H2O}
(from \ce{OH-}, now neutral) and \ce{OH-} (the oxygen of the former water
keeps the pair, charge $-1$). For ammonia: an arrow from the N \omterm{def:b1:lewis-resonance-vsepr:lewis}{lone pair} to
an H of \ce{H3O+}, and one from that O--H bond to the oxygen: N becomes
$+1$, O becomes 0.
\end{solution}

\begin{solution}{exo:b1:electronic-effects:5}
Ethanol (15.9) $<$ water (14.0) $<$ phenol (9.99, ring delocalisation) $<$
4-nitrophenol (7.16, \ce{NO2} takes the charge) $<$ ethanoic acid (4.76,
two equivalent oxygens) $<$ trichloroethanoic acid (0.51, three $-I$
chlorines).
\end{solution}

\begin{solution}{exo:b1:electronic-effects:6}
Aniline (4.60) $<$ pyridine (5.23) $<$ ammonia (9.25) $<$ methylamine
(10.66). In aniline the \omterm{def:b1:lewis-resonance-vsepr:lewis}{lone pair} is shared with the ring
(structures with C=N$^+$ and the negative charge \textit{ortho} and
\textit{para}); protonating N would cost that delocalisation.
\end{solution}

\begin{solution}{exo:b1:electronic-effects:7}
In 3-nitrophenoxide, the \omterm{def:b1:lewis-resonance-vsepr:resonance}{resonance structures} place the charge on the
carbons \textit{ortho} and \textit{para} to the oxygen, never on the carbon
bearing \ce{NO2}: no mesomeric help, hence weaker than the 4-isomer. The
nitro group still pulls electrons through the $\sigma$ bonds ($-I$),
hence stronger than phenol.
\end{solution}

\begin{solution}{exo:b1:electronic-effects:8}
\ce{CH3+} $<$ \ce{CH3CH2+} $<$ \ce{CH2=CH-CH2+} $<$ \ce{C6H5CH2+}
$\approx$ \ce{(CH3)3C+} $<$ \ce{CH3OCH2+}: alkyl groups give electrons;
allyl and benzyl delocalise the charge (two and four structures); the
oxygen \omterm{def:b1:lewis-resonance-vsepr:lewis}{lone pair} gives a structure with a full octet everywhere. The order
in the middle of the list depends on the medium.
\end{solution}

\begin{solution}{exo:b1:electronic-effects:9}
\ce{OH-} (\omterm{def:b1:electronic-effects:nucleophile}{nucleophile}) attacks the carbon of \ce{CH3Br} (\omterm{def:b1:electronic-effects:nucleophile}{electrophile}),
the C--Br pair leaving on Br. \ce{H2O} gives a pair to \ce{H+} (here acid
and base are the names used). \ce{CN-} attacks the carbonyl carbon of
ethanal, the C=O $\pi$ pair moving to oxygen. \ce{NH3} gives its \omterm{def:b1:lewis-resonance-vsepr:lewis}{lone pair}
to the empty orbital of boron in \ce{BF3}.
\end{solution}

\begin{solution}{exo:b1:electronic-effects:10}
No: $K_a = 10^{-0.51} = 0.31$ is much larger than $C$; the acid is almost
totally dissociated. $x^2/(0.010 - x) = 0.31$ gives $x = \qty{9.7e-3}{mol/L}$,
$\mathrm{pH} = 2.01$, \omterm{def:b1:predominant-reaction:dissociation}{degree of dissociation} \qty{97}{\percent}.
\end{solution}

\begin{solution}{exo:b1:electronic-effects:11}
Alkyl groups push electrons ($+I$) onto the alkoxide oxygen and make its
charge less stable, the more so the more groups. Larger alkoxides are
also less well solvated by water, which stabilises a small charged
oxygen better than one buried among alkyl groups. Both raise the
$\mathrm{p}K_a$.
\end{solution}

\begin{solution}{exo:b1:electronic-effects:12}
In an amide the nitrogen \omterm{def:b1:lewis-resonance-vsepr:lewis}{lone pair} is shared with the C=O
(\ce{N+=C-O^-} structure); taking it for a proton or an \omterm{def:b1:electronic-effects:nucleophile}{electrophile} would
destroy that stabilisation. In \omterm{def:b1:predominant-reaction:strong-acid}{strong acid} the amide is protonated on the
oxygen, which carries the partial negative charge.
\end{solution}

\begin{solution}{pb:b1:electronic-effects:1}
\textbf{1.} $10^{1.89} = 78$; $10^{1.52} = 33$; $10^{0.84} = 6.9$.
\textbf{2.} No: each further chlorine helps less. The charge of the
carboxylate is already well spread, and the measurements become less
precise as the acid approaches total dissociation.
\textbf{3.} Each Cl attracts electron density through the C--C and C--O
bonds and stabilises the carboxylate ($-I$).
\textbf{4.} The \omterm{def:b1:electronic-effects:inductive}{inductive effect} fades within two or three bonds.
\textbf{5.} The two acids appear equally strong (0.52 and 0.51). In water,
both are almost totally dissociated at usual concentrations
(\cref{ch:b1:predominant-reaction}, levelling): a $\mathrm{p}K_a$ near 0
is the limit of what water can distinguish, and the two values carry
uncertainties of a few tenths.
\textbf{6.} Ethanoic acid: acid form (\qty{98}{\percent}); chloroethanoic:
both, the anion slightly predominant ($10^{3 - 2.87} = 1.3$); the three
others: anion.
\textbf{7.} \ce{CH3-C(=O)-O^-} $\leftrightarrow$ \ce{CH3-C(-O^-)=O}.
\textbf{8.} Two equal C--O bonds, between a double and a single bond.
\textbf{9.} The charge is on the single oxygen, with no $\pi$ bond to share
it.
\textbf{10.} $10^{15.9 - 4.76} = 10^{11.1}$, about $10^{11}$.
\textbf{11.} The phenoxide charge spreads to three ring carbons, less
electronegative than oxygen: a gain, smaller than in a carboxylate:
$\mathrm{p}K_a$ 9.99 between 4.76 and 15.9.
\textbf{12.} $K = 10^{6.37 - \mathrm{p}K_a}$: ethanoic acid $10^{1.61} = 41$,
favoured (carbon dioxide is released with excess hydrogencarbonate);
phenol $10^{-3.62}$ and ethanol $10^{-9.5}$: no reaction.
\textbf{13.} 4-Nitrophenol (7.16) $>$ 2-nitrophenol (7.23) $>$ 3-nitrophenol
(8.36) $>$ phenol (9.99).
\textbf{14.} The quinoid structure of \cref{exo:b1:electronic-effects:3},
with the charge on an oxygen of the nitro group.
\textbf{15.} The charge reaches only the carbons \textit{ortho} and
\textit{para} to the oxygen; the nitro carbon is \textit{meta}.
\textbf{16.} The $-I$ effect of \ce{NO2}.
\textbf{17.} In 2-nitrophenol the OH forms a \omterm{def:b1:intermolecular-forces-solvents:hbond}{hydrogen bond} with an oxygen
of the neighbouring nitro group, which stabilises the acid form and makes
the proton slightly harder to remove.
\textbf{18.} Anion fractions $1/(1 + 10^{\mathrm{p}K_a - 7.5})$: 4-nitro
\qty{69}{\percent}, 2-nitro \qty{65}{\percent}, 3-nitro
\qty{12}{\percent}. Acid and anion of 4-nitrophenol absorb differently,
the anion in the visible: its colour appears around its $\mathrm{p}K_a$.
\textbf{19.} $x^2/(0.010 - x) = 10^{-0.52} = 0.30$: $x = \qty{9.7e-3}{mol/L}$,
$\mathrm{pH} = 2.01$: almost a \omterm{def:b1:predominant-reaction:strong-acid}{strong acid}.
\textbf{20.} $\frac12(4.76 + 2.00) = 3.38$.
\textbf{21.} \ce{CF3COOH + CH3COO- -> CF3COO- + CH3COOH}, $K = 10^{4.76 -
0.52} = 10^{4.24}$: \omterm{def:b1:extent-q-and-k:quantitative}{quantitative}.
\textbf{22.} $K_a(\ce{CF3COOH})/K_a(\ce{CH3COOH}) =$ \textbf{$10^{4.24} =
\num{1.7e4}$}, about ten thousand.
\end{solution}
